\documentclass[10pt,a4paper]{amsart}
\usepackage[margin=19mm]{geometry}
\usepackage{amsmath,amssymb,mathtools}
\usepackage[scr=rsfs,cal=euler]{mathalfa}
\usepackage{xfrac}
\usepackage[unicode,hidelinks,bookmarks=false]{hyperref}
\newcommand{\ie}{i.e.\ }
\newcommand{\cf}{cf.\ }
\newcommand{\resp}{respectively\ }
\newcommand{\Subsection}{\subsection}
\providecommand{\nobreakdash}{\nobreak\hbox{-}}
\setlength{\emergencystretch}{2em}
\multlinegap0pt
\usepackage{amssymb}
\usepackage{mathtools}
\usepackage{tikz}
\usepackage[shortlabels]{enumitem}
\newtheorem{theorem}{Theorem}
\newcommand{\F}{{\mathcal{F}}}
\newcommand{\N}{\mathbb{N}}
\newcommand{\U}{U}
\newcommand{\UU}{{\mathcal{W}}}
\newcommand{\X}{{\mathbb{X}}}
\newcommand{\cdim}{\mathrm{cdim}}
\renewcommand{\dim}{{\mathrm{dim}}}
\newcommand{\x}{x}	
\title{Point-to-set Principle and Constructive Dimension Faithfulness}
\author{Satyadev Nandakumar, Subin Pulari, Akhil S}
\date{Source: arXiv:2403.08278v4 (CC BY 4.0)}
\begin{document}
\maketitle

\noindent\emph{Source and licence.} Point-to-set Principle and Constructive Dimension Faithfulness. Version arXiv:2403.08278v4 (CC BY 4.0), distributed by arXiv under the Creative Commons Attribution 4.0 International licence, http://creativecommons.org/licenses/by/4.0/, as recorded in the arXiv licence metadata for this exact version and shown on the abstract page https://arxiv.org/abs/2403.08278v4. Full licence notice: \texttt{licenses/arxiv-2403.08278-CC-BY-4.0.txt}.\par\smallskip

\noindent\textit{Editorial scope.} Counted unit: the article's theorem thm:pspforPhidim (source0.tex lines 896--899). The article's complete original proof of it is delivered with it (source0.tex lines 900--930, one \texttt{proof} environment of 31 lines). No mathematics is written, repaired or re-derived: the statement and the proof are the article's own lines, in the article's own notation, and no retained line is substituted or re-flowed.  No further source-local context is needed: every cross-reference of every retained line resolves inside the two delivered spans.  Nothing was omitted from any retained span, so the package carries no omission record.  No retained line contains a citation, so the wrapper carries no bibliography and no bibliography entry.  No retained line contains a figure environment, an image-inclusion command or drawing code, so no image is delivered and the package declares no asset dependency.  Editorial formatting only, in addition: an AMS \texttt{amsart} preamble that restates the article's own declarations and macros used by the retained lines, namely the environments \texttt{theorem} on separate counters, together with the standard packages those lines need.  The retained environments print in this document as Theorem 1 in order of appearance, whereas the article prints the same items with the numbering of its own full document.  The complete native source of the article as distributed in the arXiv e-print for this exact version is bundled unchanged as \texttt{sources/156013/source0.tex}.
\begin{theorem}\label{thm:pspforPhidim}
	Let $\Phi$ be a family of  covering sets over $\X$. For all $\mathcal{F} \subseteq \X$, 
	$$\dim_\Phi(\mathcal{F}) = \min\limits_{A \subseteq \N} \; \sup\limits_{\x \in \mathcal{F}} \; \liminf_{r \rightarrow \infty} \frac{K_r^A(\x, \Phi)}{r}.$$
\end{theorem}

\begin{proof}
	
	By the definition of $\Phi$-dimension, for any rational $s > \dim_\Phi(\mathcal{F})$ and $k \in \mathbb{N}$, there exists a sequence of $\Phi$-covers $\{\U_i^{k,s}\}_{i \in \mathbb{N}}$ of $\mathcal{F}$ such that
	\(
	\sum_{i \in \mathbb{N}} |\U_i^{k,s}|^s \leq 2^{-k}.
	\)
	
	 Consider the oracle sequence $A \subseteq \mathbb{N}$ that, for each $i, k \in \mathbb{N}$ and rational $s > \dim_\Phi(\mathcal{F})$, encodes:
	\begin{itemize}
		\item $f(i,k,s) = \delta^{-1}(\U_i^{k,s})$, the index of the $i$-th cover set, and
		\item $g(i,k,s) = 2^{-r}$, where $r \in \mathbb{N}$ is the unique integer satisfying $2^{-r-1} < |\U_i^{k,s}| \leq 2^{-r}$.
	\end{itemize}

	We show that given oracle access to $A$, for any rational $s > \dim_\Phi(\mathcal{F})$ and	for any point $x \in \F$, ${K_r^A(\x, \Phi)} \leq rs + o(r)$ for infinitely many $r \in \N$.
	
	Given $m \in \mathbb{N}$, choose $k > ms$. Since $\sum_{i \in \mathbb{N}} |\U_i^{k,s}|^s \leq 2^{-k} \leq 2^{-ms}$, every set $U$ in the cover satisfies $|U| < 2^{-m}$.
		
	
	For any $x \in \F$, consider the following description of the set in $\{\U_i^{k,s}\}_{i \in \mathbb{N}}$ that contains $x$. We specify the diameter range (an $r$ such that $2^{-r-1} < |U| \leq 2^{-r}$) of the set, and an index $j$ within sets of that diameter range. 
	Since $\sum_{i \in \mathbb{N}} |\U_i^{k,s}|^s \leq 1$, for any fixed $r \in \mathbb{N}$ there are at most $2^{rs + s} $ sets $U$ in the cover satisfying $2^{-r-1} < |U|$. Hence each such set can be indexed using $rs + o(r)$ bits. 
	
	More precisely, consider the Turing machine $M$ with oracle access to $A$ that, on input $(r, j, s, k)$, proceeds as follows: iterate through all $i \in \mathbb{N}$ and use $A$ to check whether $g(i,k,s) = r$. Upon finding the $j$-th such index $i$, halt and output $f(i,k,s)$.
	The input $(r, j, s, k)$ is encoded using $rs + o(r)$ bits to specify $j$ and $o(r)$ bits to specify $r,s,k$. Therefore, for any $m \in \N$, we have that for some $r \geq m$, $K_r^A(x, \Phi) \leq rs + o(r)$.
	
	\medskip
	
	We now show the reverse direction of the inequality.
	For an $A \in \Sigma^\infty$, take a rational $s > \sup\limits_{\x \in \mathcal{F}} \; \liminf\limits_{r \rightarrow \infty} \frac{K_r^A(\x, \Phi)}{r}$. For all $x \in \F$, we have that for some $r \in \N$, there exists a $U \in \Phi$, with $|U| < 2^{-r}$ and $K^A(\delta^{-1}(U)) \leq sr$. Define a $\Phi$-cover $\UU$ as follows. For each $w \in \Sigma^*$, if $K^A(w) \leq rs$ for some $r \in \N$ such that $ |\delta(w)| \leq 2^{-r}$, then add $\delta(w)$ to $\UU$.
	It follows that $\F \subseteq \bigcup_{U \in \UU} U $.
	We also have $\sum_{U \in \UU_k} |U|^s \leq \sum_{w \in \Sigma^*} 2^{-K(w)} \leq 1$. The last inequality follows from the Kraft inequality. Therefore $\cdim_{\Phi}(\F) \leq s$.
\end{proof}

\end{document}
